% META: {"id":"1SPE-PROBCOND-EX-028","chapitre":"1SPE-PROBA-COND","type_objet":"exercice","capacites":["1SPE-PROBA-COND-C3"],"capacites_codes":["C3"],"methodes":["M3"],"parcours":2,"competences":["calculer","raisonner"],"duree_min":15,"mode_creation":"ex_nihilo","sources_inspiration":[],"parametres_sympy":{},"fichier_tex":"chapitres/1SPE-PROBA-COND/exercices/1SPE-PROBCOND-EX-028.tex","corrige_tex":"chapitres/1SPE-PROBA-COND/corriges/1SPE-PROBCOND-CO-028.tex","status":"approved","origin":"generated","status_history":["generated","audited","accepted","approved"],"release_acceptance":"RELEASE_OWNER_BATCH_ACCEPTANCE","acceptance_closure_digest":"sha256:9b3ccf9a81c5520fbb7e03b7b2d3e7bf2057a02834b6d908bf3be5f49f3b553f","acceptance_receipt_digest":"sha256:4fad86f0282e0fa4e0419cca1ad6379ae7af5a280c04a4a90880f90a1c044242"}
% BEGIN-VERIFY
% from sympy import Rational
% # Rediger la demonstration de la formule des prob totales
% # Verification numerique
% P_B = Rational(3, 8)
% P_A_B = Rational(2, 3)
% P_A_Bbar = Rational(1, 4)
% P_A = P_B * P_A_B + (1-P_B) * P_A_Bbar
% assert P_A == Rational(3,8)*Rational(2,3) + Rational(5,8)*Rational(1,4)
% assert P_A == Rational(1,4) + Rational(5,32)
% assert P_A == Rational(8,32) + Rational(5,32)
% assert P_A == Rational(13, 32)
% END-VERIFY
\begin{exercice}{1SPE-PROBCOND-EX-028}{2}{15}
\begin{enumerate}
  \item Rediger la démonstration de la formule des probabilités totales dans le cas de la partition $\{B, \overline{B}\}$.

  \item \textbf{Application.} On donne $P(B) = \dfrac{3}{8}$, $P_B(A) = \dfrac{2}{3}$ et $P_{\overline{B}}(A) = \dfrac{1}{4}$. Calculer $P(A)$.
\end{enumerate}
\end{exercice}
